\chapter*{Remerciements}

La r\'ealisation de ce projet de fin d'\'etudes s'est appuy\'ee sur l'accompagnement et le soutien de plusieurs personnes auxquelles une sinc\`ere gratitude est adress\'ee.

Des remerciements particuliers sont d'abord adress\'es \`a Monsieur \textbf{Badis Messaoudi}, encadrant au sein de l'entreprise d'accueil, pour la qualit\'e de son suivi, la pertinence de ses orientations et la disponibilit\'e accord\'ee tout au long du stage. Son accompagnement a permis d'aborder le projet \textit{HireCue} dans un cadre professionnel exigeant, avec un niveau d'attention constant port\'e \`a la coh\'erence technique et \`a la valeur metier des \'evolutions r\'ealis\'ees.

Une reconnaissance particuli\`ere est \'egalement adress\'ee \`a Monsieur \textbf{Chedy Bouslema}, encadrant acad\'emique, pour ses conseils, ses remarques constructives et son appui m\'ethodologique. Ses orientations ont contribu\'e \`a mieux structurer ce travail et \`a maintenir une lecture acad\'emique coh\'erente des choix techniques et fonctionnels pr\'esent\'es dans le rapport.

Des remerciements sont aussi adress\'es \`a l'ensemble de l'\'equipe de travail pour l'accueil, l'esprit de collaboration et les \'echanges techniques qui ont facilit\'e l'appropriation du projet et de son environnement. Le cadre de travail offert pendant le stage a constitu\'e un contexte favorable \`a l'apprentissage, \`a l'approfondissement des comp\'etences mobilis\'ees et \`a une meilleure compr\'ehension des contraintes r\'eelles d'une plateforme comme \textit{HireCue}.

Une pens\'ee respectueuse est \'egalement adress\'ee aux membres du jury pour l'attention port\'ee \`a ce rapport et pour le temps consacr\'e \`a son \'evaluation.

Enfin, une gratitude sinc\`ere est exprim\'ee \`a toutes les personnes qui ont apport\'e, directement ou indirectement, un soutien intellectuel, moral ou organisationnel au cours de cette exp\'erience.
